% Research contribution for integration into the parent article.
% Requires amsmath, amssymb, amsthm, booktabs, and theorem/lemma/corollary environments.
% All labels use the prefix fg:.
\section{The full Gaussian single-product rank conjecture}
\label{fg:sec:rank}

The manuscript's Conjecture~\texttt{signed:conj:rank} concerns a precise
coefficient matrix, rather than the dimension of a space of actual
polylogarithmic numbers. We prove both proposed rank formulas in every
odd weight. The proof also gives the full even-weight answer, an explicit
basis of the nullspace, and a complete finite test for whether an arbitrary
coefficient vector follows from the stated rows.

The additional boundary issue is essential. A preceding incoming report,
\path{polylogarithms_research_2026-10-10.zip}, proves the dimension after
setting the Gaussian family to zero and after retaining regularized rows
in a different presentation. Our result keeps that family free, removes
exactly the manuscript's divergent coordinate, and retains only the
stuffle-minus-shuffle difference when a single factor diverges. The new
step is a canonical lift of every allowed Gaussian nullvector through
that precise boundary presentation.

\subsection{The matrix and its scope}

Fix an integer $w\ge2$ and write $n=w-2$. For $a+b=w$, let
\[
 D_{a,b}(x,y)=\sum_{k>\ell>0}\frac{x^ky^\ell}{k^a\ell^b},
 \qquad X_a^{r,s}=\operatorname{Im}D_{a,w-a}(i^r,i^s),
 \quad r,s\in\mathbb Z/4\mathbb Z.
\]
For the coefficient computation, the $X_a^{r,s}$ are formal rational
coordinates. Impose conjugation
$X_a^{-r,-s}=-X_a^{r,s}$ and put $X_a^{r,s}=0$ if both colors are even.
A set of representatives of the remaining color pairs is
\begin{equation}\label{fg:eq:pairs}
 (0,1),\ (1,0),\ (1,1),\ (1,2),\ (1,3),\ (2,1).
\end{equation}
Remove $X_1^{0,1}$, the single divergent imaginary coordinate in these
representatives. Thus the ambient coefficient space has dimension
$N_w=6(w-1)-1=6w-7$.

For $p+q=w$ and $r,s\in\mathbb Z/4\mathbb Z$, define the stuffle and
shuffle coefficient rows
\begin{align}
 S_{p,q}^{r,s}
 &=X_p^{r,s}+X_q^{s,r},\\
 H_{p,q}^{r,s}
 &=\sum_{j=0}^{p-1}\binom{q+j-1}{j}X_{q+j}^{s,r-s}
   +\sum_{j=0}^{q-1}\binom{p+j-1}{j}X_{p+j}^{r,s-r}.
 \label{fg:eq:raw-rows}
\end{align}
If both single factors $\operatorname{Li}_p(i^r)$ and
$\operatorname{Li}_q(i^s)$ are finite, retain the two rows separately.
If exactly one factor is $\operatorname{Li}_1(1)$, retain only
$S_{p,q}^{r,s}-H_{p,q}^{r,s}$, in which the removed coordinate cancels.
The two-divergent-factor case occurs only at $w=2$ with two real colors;
its imaginary coefficient row is zero. Discard zero rows; retaining
repeated rows does not change the problem. Call the resulting rational
matrix $A_w$. Let $A_w^{\neg g}$ be obtained by deleting the $w-1$
columns $X_a^{1,0}$.

The depth-one right sides are suppressed because the present question
concerns the coefficient row space. Analytically, these rows come from
convergent shuffle and stuffle identities, and from subtracting the
one-divergence identities before taking a regularized limit. Nothing
here assigns a value to $\zeta(1)$. This is a restricted linear system
of products of two depth-one values. In particular, it omits distribution,
higher-depth products and elimination, and additional argument
transformations. It is narrower than the families of standard relations
studied by Zhao; see J.~Zhao, \emph{Standard Relations of Multiple
Polylogarithm Values at Roots of Unity}, Documenta Mathematica
\textbf{15} (2010), 1--34, \texttt{arXiv:0707.1459}.

\begin{theorem}[Exact ranks in every weight]\label{fg:thm:rank}
The matrix family just defined satisfies
\begin{align}
 \operatorname{rank}_{\mathbb Q}A_w
 &=\begin{cases}
      6w-7,&w\text{ even},\\
      5w-6,&w\text{ odd},
    \end{cases}\label{fg:eq:full-rank}\\
 \operatorname{rank}_{\mathbb Q}A_w^{\neg g}
 &=\begin{cases}
      5w-6,&w\text{ even},\\
      (9w-11)/2,&w\text{ odd}.
    \end{cases}\label{fg:eq:nong-rank}
\end{align}
In particular, both equalities in the manuscript's odd-weight rank
conjecture hold for every odd $w\ge3$.
\end{theorem}

The proof is an explicit description of the entire kernel. It does not
extrapolate from finite rank calculations.

\subsection{Encoding the rows by homogeneous polynomials}

For a scalar assignment to the formal coordinates, set
\[
 \mathscr F^{r,s}(X,Y)
   =\sum_{a=1}^{w-1}X_a^{r,s}X^{a-1}Y^{w-a-1}.
\]
These homogeneous polynomials have degree $n$. Write
\[
 U=\mathscr F^{0,1},\quad G=\mathscr F^{1,0},\quad
 A=\mathscr F^{1,1},\quad K=\mathscr F^{1,2},\quad
 M=\mathscr F^{1,3},\quad V=\mathscr F^{2,1}.
\]
The letter $A$ in these formulas denotes a polynomial; the matrix is
always written $A_w$. The removed coordinate is encoded by
\begin{equation}\label{fg:eq:removed}
 [Y^n]U(X,Y)=0.
\end{equation}
For a polynomial $P$, let $P^\sigma(X,Y)=P(Y,X)$.

Coefficient comparison turns a complete set of finite stuffle rows into
$\mathscr F^{r,s}(X,Y)+\mathscr F^{s,r}(Y,X)=0$. Similarly, the shuffle
rows are the coefficients of
\[
 \mathscr F^{s,r-s}(X+Y,X)
   +\mathscr F^{r,s-r}(X+Y,Y).
\]
For example, the coefficient of $X^{p-1}Y^{q-1}$ in the second summand
is precisely the second binomial sum in
\eqref{fg:eq:raw-rows}. This explains the substitutions without invoking
any rank or transcendence theorem.

\begin{lemma}[The boundary defect]\label{fg:lem:boundary}
Put $c=[X^n]G(X,Y)$. The finite rows with product colors $(0,1)$,
together with their one-divergence difference row, are exactly
\begin{align}
 U(X,Y)+G(Y,X)&=cY^n,\\
 M(X+Y,X)+U(X+Y,Y)&=cY^n.
 \label{fg:eq:boundary}
\end{align}
Consequently they determine
\begin{equation}\label{fg:eq:UM}
 U(X,Y)=-G(Y,X)+cY^n,
 \qquad M(X,Y)=G(X-Y,X).
\end{equation}
\end{lemma}
\begin{proof}
For colors $(0,1)$, the factor $\operatorname{Li}_p(1)$ is finite
exactly when $p\ge2$. Therefore the coefficients of
$X^{p-1}Y^{q-1}$ in the stuffle polynomial
$U(X,Y)+G(Y,X)$ vanish for all $p\ge2$. The only possible defect is
a multiple of $Y^n$. By \eqref{fg:eq:removed}, its coefficient is $c$.
The same finite-product argument shows that
$M(X+Y,X)+U(X+Y,Y)$ has only a $Y^n$ defect, say $dY^n$.
When $p=1$, the retained regularized row is the difference of these
two coefficient rows. It says exactly $c-d=0$. This proves
\eqref{fg:eq:boundary}. Substituting the first identity into the
second cancels $cY^n$ and gives
$M(X+Y,X)=G(Y,X+Y)$. The invertible substitution
$(X+Y,X)\mapsto(X,Y)$ yields \eqref{fg:eq:UM}.
The argument also covers $w=2$: there are no $p\ge2$ coefficients,
and the difference row still equates the two constant defects.
\end{proof}

This lemma is the point at which the precise regularization convention
matters. Setting the two boundary rows separately to zero would impose
an extra condition on $c$. Deleting the boundary row altogether would
lose a condition. The manuscript's difference convention produces the
unique correction $cY^n$.

\begin{lemma}[Exhaustive kernel conditions]\label{fg:lem:conditions}
An assignment annihilates every row of $A_w$ if and only if its six
polynomials satisfy \eqref{fg:eq:UM} and
\begin{align}
 A(Y,X)&=-A(X,Y), & V(X,Y)&=-K(Y,X),\\
 M(Y,X)&=M(X,Y), & G(X,Y)&=-G(X,X-Y),\\
 K(X,Y)&=-A(Y,Y-X),&K(X,Y)&=K(X,X-Y).
 \label{fg:eq:conditions}
\end{align}
\end{lemma}
\begin{proof}
The product-color classes $(1,1),(1,2),(1,3)$ give the conditions in
this table:
\[
\begin{array}{c|c|c}
 \text{colors}&\text{stuffle}&\text{shuffle}\\ \hline
 (1,1)& A+A^\sigma=0
   &G(X+Y,X)+G(X+Y,Y)=0\\
 (1,2)&K+V^\sigma=0
   &A(X+Y,X)-V(X+Y,Y)=0\\
 (1,3)&M-M^\sigma=0
   &K(X+Y,Y)-K(X+Y,X)=0.
\end{array}
\]
For the $(1,2)$ shuffle we have exchanged the two product factors;
this merely transposes $X,Y$ and retains exactly the same collection
of coefficient rows. Substituting $V=-K^\sigma$ into its shuffle gives
$K(X,Y)=-A(Y,Y-X)$. The other two shuffle equations yield the two
reflection equations in \eqref{fg:eq:conditions}.

There are no omitted nonzero color classes. Up to swapping the product
factors and conjugation, all unordered pairs are
\[
 (0,0),\ (0,1),\ (0,2),\ (1,1),\ (1,2),\ (1,3),\ (2,2).
\]
The classes $(0,0),(0,2),(2,2)$ are entirely real and their imaginary
rows vanish. The $(0,1)$ class is exactly
Lemma~\ref{fg:lem:boundary}. This proves necessity and sufficiency.
\end{proof}

\subsection{The two independent polynomial sectors}

\begin{theorem}[Canonical description of the full nullspace]
\label{fg:thm:kernel}
For even $w$, $\ker A_w=0$. For odd $w=2m+1$, its elements are in
one-to-one correspondence with independent choices of degree
$n=2m-1$ homogeneous polynomials $G,A$ satisfying
\begin{equation}\label{fg:eq:two-sectors}
 G(X,Y)=-G(X,X-Y),\qquad A(X,Y)=-A(Y,X).
\end{equation}
All six polynomials are then given by
\begin{align}
 U(X,Y)&=-G(Y,X)+cY^n,&c&=[X^n]G,\\
 M(X,Y)&=G(X-Y,X),\\
 K(X,Y)&=-A(Y,Y-X),\\
 V(X,Y)&=-K(Y,X).
 \label{fg:eq:kernel-map}
\end{align}
Each sector has dimension $m$. In particular,
$\dim\ker A_w=2m=w-1$ at odd weight.
\end{theorem}
\begin{proof}
The preceding lemmas show that every kernel element is of the form
\eqref{fg:eq:kernel-map}, with the two conditions
\eqref{fg:eq:two-sectors}. It remains to impose the symmetry of $M$
and the reflection symmetry of $K$. Homogeneity and the reflection
condition on $G$ give
\begin{align*}
 M(Y,X)&=G(Y-X,Y)
       =(-1)^nG(X-Y,-Y)\\
       &=-(-1)^nG(X-Y,X)
        =(-1)^{n+1}M(X,Y).
\end{align*}
Likewise, antisymmetry of $A$ gives
\begin{align*}
 K(X,X-Y)&=-A(X-Y,-Y)\\
         &=-(-1)^nA(Y-X,Y)
          =(-1)^nA(Y,Y-X)\\
         &=(-1)^{n+1}K(X,Y).
\end{align*}
If $n$ is even, the required equalities $M=M^\sigma$ and
$K(X,Y)=K(X,X-Y)$ force $M=K=0$, hence $G=A=0$, and then all six
polynomials vanish. If $n$ is odd, these two remaining equalities
hold automatically. The exhaustive list of classes in
Lemma~\ref{fg:lem:conditions} proves that every pair in
\eqref{fg:eq:two-sectors} is a kernel element, so this is a bijective
parameterization.

For the $G$ sector, change variables from $X,Y$ to $X,Z=2Y-X$.
Reflection changes $Z$ to $-Z$, so $G$ is an odd polynomial in $Z$.
At total degree $2m-1$, a basis is
\begin{equation}\label{fg:eq:G-basis}
 G_j(X,Y)=X^{2j}(2Y-X)^{2m-1-2j},\qquad 0\le j<m.
\end{equation}
For the $A$ sector, a basis is
\begin{equation}\label{fg:eq:A-basis}
 A_j(X,Y)=X^{2m-1-j}Y^j-X^jY^{2m-1-j},\qquad 0\le j<m.
\end{equation}
Their independence follows respectively from the distinct powers of
$Z$, and from their disjoint exchanged monomial pairs. Thus each
sector has dimension $m$.
\end{proof}

\begin{proof}[Proof of Theorem~\ref{fg:thm:rank}]
The full matrix has $6w-7$ columns. Subtracting the kernel dimensions
in Theorem~\ref{fg:thm:kernel} gives
\eqref{fg:eq:full-rank}. To pass to $A_w^{\neg g}$, extend a vector
on the remaining columns by zero on the deleted Gaussian columns.
It is in $\ker A_w^{\neg g}$ if and only if the extended vector is
in $\ker A_w$ and $G=0$. At odd weight this leaves exactly the
$m$-dimensional $A$ sector; at even weight it leaves zero. There
are $(6w-7)-(w-1)=5w-6$ remaining columns. Subtracting these
kernel dimensions gives \eqref{fg:eq:nong-rank}.
\end{proof}

The Hilbert series of these formal nullspaces are especially simple:
\[
 \sum_{w\ge2}\dim\ker A_w\,t^w
 =\frac{2t^3}{(1-t^2)^2},\qquad
 \sum_{w\ge2}\dim\ker A_w^{\neg g}\,t^w
 =\frac{t^3}{(1-t^2)^2}.
\]
The second is the previously obtained quotient count; the first records
the additional freely liftable Gaussian sector.

\subsection{Every Gaussian-only consequence is already a same-point shuffle}

Let $\mathcal R_w$ be the row space of $A_w$ and let $E_g$ be the
coordinate subspace supported only on the Gaussian columns.
Projection of $\mathcal R_w$ onto the other columns has dimension
$\operatorname{rank}A_w^{\neg g}$. Hence
\begin{equation}\label{fg:eq:intersection}
 \dim(\mathcal R_w\cap E_g)
 =\operatorname{rank}A_w-\operatorname{rank}A_w^{\neg g}
 =\begin{cases}
    w-1,&w\text{ even},\\
    (w-1)/2,&w\text{ odd}.
  \end{cases}
\end{equation}

\begin{corollary}[Completeness for Gaussian-only coefficient rows]
\label{fg:cor:g-only}
For odd $w=2m+1$, every rational linear combination of the defining
rows whose non-Gaussian coefficients cancel is a rational linear
combination of the $m$ same-point product-shuffle rows obtained from
\[
 \operatorname{Li}_p(i)\operatorname{Li}_{w-p}(i),
 \qquad 1\le p\le m.
\]
Thus enlarging a search to all the colors and one-divergence rows in
$A_w$ yields no additional Gaussian-only coefficient relation in
odd weight beyond these elementary shuffles.
\end{corollary}
\begin{proof}
These $m$ rows are present in $A_w$ and are supported only on $G$.
They are independent: in the usual inner-index ordering, the row
with $p\le m$ has largest inner index $q=w-p$, with coefficient
one, and no larger inner index occurs. Ordering the selected rows
and columns accordingly gives a triangular square block with unit
diagonal. Their span is therefore an $m$-dimensional subspace of
$\mathcal R_w\cap E_g$, which also has dimension $m$ by
\eqref{fg:eq:intersection}. The spaces are equal.
\end{proof}

This is a statement about the coefficient projection of the equations.
When actual depth-one right sides are restored, any equality of row
coefficients also gives the corresponding identity of those right
sides, because all the underlying analytic rows are valid. The
corollary does not forbid Gaussian-only identities from other
functional equations or higher-depth elimination.

\subsection{A complete membership test and exact certificates}

Equations \eqref{fg:eq:G-basis}--\eqref{fg:eq:A-basis}, followed by
\eqref{fg:eq:kernel-map}, give $w-1$ explicit integer vectors
$v_1,\ldots,v_{w-1}$ in odd weight. For any rational coefficient row
$t$, elementary finite-dimensional duality gives
\begin{equation}\label{fg:eq:membership}
 t\in\operatorname{rowspan}_{\mathbb Q}A_w
 \quad\Longleftrightarrow\quad
 t\cdot v_j=0\quad\text{for every }j.
\end{equation}
Thus membership in the stated coefficient module has a complete test
using only integer polynomial substitution and rational dot products.
A nonzero dot product supplies a short exact separating functional.
At even weight the row space is the whole coordinate space. None of
these tests asserts the absence of a relation outside this matrix.

At weight five, $n=3$, one may choose
\[
 G_0=(2Y-X)^3,\quad G_1=X^2(2Y-X),\qquad
 A_0=X^3-Y^3,\quad A_1=X^2Y-XY^2.
\]
These four parameters give the four-dimensional nullspace of the
$92\times23$ matrix, with no numerical polylogarithm computation.
The $G$-zero subspace has the two $A$ parameters, explaining the
rank pair $(19,17)$. At weight seven the corresponding pair is
$(29,26)$, and at weight thirty-one it is $(149,134)$.

The supplied standard-library program
\texttt{full\_gaussian\_rank.py} constructs the raw rows from
\eqref{fg:eq:raw-rows}; its separate polynomial code constructs
\eqref{fg:eq:kernel-map}. It verifies that every displayed basis
vector annihilates every raw row. Independent rational elimination
checks both ranks in weights $2$ through $11$. Basis identities are
also replayed in weights $12$ through $17$ and in weights
$23,31,47,63$. The JSON receipt includes the entire weight-five
matrix, coordinate order, row labels, and kernel basis. These finite
checks validate the implementation; the preceding proof establishes
all weights.

\subsection{Integration notes and further questions}

The manuscript's Conjecture~\texttt{signed:conj:rank} can be replaced
by Theorem~\ref{fg:thm:rank}. The even-weight formulas and the
canonical kernel description are additional results. The word
``equivalently'' preceding the manuscript's row-space-intersection
formulation should be changed to ``consequently,'' unless one of
the absolute rank formulas is taken as an additional premise: a
specified difference of two ranks does not by itself determine both
ranks. Our proof establishes the absolute ranks and their difference.

Three directions now have precise formulations. First, replace the
implicit numerical rank routine by the explicit polynomial test
\eqref{fg:eq:membership} in future relation searches. Second, add
specified distribution or octahedral rows and compute their action
on these two small polynomial sectors, rather than repeating
elimination in all $6w-7$ coordinates. Third, classify which
higher-depth relations couple the two sectors. Such a classification
would address a larger relation module; passing from any formal
module to numerical period independence remains a separate problem.
